\UTIResumenPagina
  {spanish}
  {Resumen Ejecutivo}
  {La inferencia de un modelo de lenguaje se divide en dos fases con cuellos de botella distintos: el prefill, que procesa el prompt, y el decode, que genera la respuesta token a token. Este trabajo se propuso caracterizar ambas fases para Qwen3-14B servido con vLLM sobre una GPU NVIDIA A100 de 40~GB del clúster HPC de CEDIA, y medir el efecto de la cuantización, la longitud de contexto y la concurrencia. Se siguió una metodología de prueba y error en once iteraciones. Ninguno de los cuatro arranques de \texttt{vllm serve}, uno con Qwen3-8B y tres con Qwen3-14B, quedó listo: la importación del entorno desde el almacenamiento compartido NFS pasó de 12,9 a 25,8~minutos entre intentos, y la copia al disco local no se completó. Por ello no se obtuvieron métricas de Qwen3-14B. Las fases se caracterizaron con Ollama sobre la A100 completa en dos modelos de 30B cuantizados en Q4\_K\_M. En \texttt{nemotron-3-nano:30b}, el tiempo por token de salida varió un 0,95~\% entre 170 y 2390 tokens de prompt, mientras la velocidad de prefill se multiplicó por 1,99; en \texttt{qwen3-coder:30b}, el primero aumentó un 7,1~\% y la segunda un 56,6~\%. Esta asimetría concuerda con un decode limitado por el ancho de banda de memoria. Se detectó, además, que en Ollama el tiempo hasta el primer token medido en el cliente incluye un \texttt{load\_duration} que llega al 94~\% de su valor, por lo que el prefill debe tomarse del servidor. Se concluye que el objetivo general no se alcanzó para la combinación prevista, y que el trabajo aporta la caracterización de la asimetría entre fases, la evidencia de la inviabilidad del arranque desde NFS en este clúster y un criterio de validez para medir con Ollama.}
  {Descriptores}
  {decode, inferencia de modelos de lenguaje, Ollama, prefill}
